\section{收束与扩展：从语言 MoE 到 Vision MoE}

讲座最后先总结 Switch 的语言建模结果，再展示视觉领域的新工作。两部分共同说明：稀疏专家的核心接口可以跨模态复用，但 token 的含义、冗余结构和重要性排序会改变 routing 策略。文本通常希望每个 token 接受类似计算；图像 patch 则可能允许更激进的选择性处理。

\lecturefigure{slide-36-wrapping-up.jpg}{讲座总结：简单路由、稳定训练、预训练加速、多语言与万亿参数规模。}{Stanford Online 官方视频 00:57:08--01:00:43。}

\begin{knowledgebox}{读图：把“结果”还原成条件句}
Switch 简化了 MoE，并在给定 T5 基线、C4/mT5 数据、TPU 与 Mesh TensorFlow 实现中取得速度和质量收益；它训练了万亿参数模型，但每 token 只激活少量 experts；它改善多语言平均表现，但仍需检查 per-language 路由与部署成本。
\end{knowledgebox}

视觉 MoE 将 image patches 当作 tokens，在 Transformer 中插入 sparse expert layers。图像存在大量局部冗余，因此 router 不仅可以决定“使用哪个 expert”，还可能决定“哪些 patches 值得占用有限 capacity”。这使 sparsity 开始接近 adaptive computation。

\lecturefigure{slide-37-sparse-models-computer-vision.jpg}{V-MoE：将 sparse expert scaling 扩展到图像分类。}{Stanford Online 官方视频 01:00:44--01:01:15。}

\begin{knowledgebox}{什么是视觉 token}
Vision Transformer 将图像切成 patches，并把每个 patch 投影为 token。V-MoE 的 expert 处理对象是 patch 表示；它不是为每张完整图片只选一个 expert。不同 patch 可被路由到不同 experts。
\end{knowledgebox}

Priority routing 在 capacity factor 小于 1 时先为 token 估计重要性，再让高优先级 patches 获得 expert 槽位。若 $\gamma=0.5$，容量只够处理理想均分的一半 token；这不再只是意外 overflow，而是有意进行计算预算分配。图像背景或重复区域可能被跳过，关键物体 patch 获得更多稀疏层计算。

\lecturefigure{slide-38-priority-routing-vision-moe.jpg}{Vision MoE 的 priority routing：在容量不足时优先处理重要 patches。}{Stanford Online 官方视频 01:01:16--01:02:10。}

\teachervoice{讲者把 capacity factor 降到 1 以下视为很有前景的方向：在图像中，甚至只让约十分之一 patch 经过 sparse layer 仍可能保持良好质量。这时 sparsity 开始同时改变“参数选择”和“计算量分配”。}

\begin{warningbox}{Priority score 也是一个学习到的假设}
被跳过的 patch 不一定真的“不重要”。若 priority router 对小物体、文字、边缘区域或少数类产生系统性低分，平均准确率可能掩盖脆弱性。应做遮挡、分群、校准与分布外测试。
\end{warningbox}

\subsection{Q\&A 给出的部署边界}

结尾 Q\&A 补充了三条 slide 上没有完整写出的边界。第一，在当时的 C4、SuperGLUE 和常见序列长度下，attention map 不是主要内存成本，权重更突出，因此线性 attention 不一定先解决瓶颈。第二，继续扩大 experts 的直接限制之一是参数必须存放在某处。第三，sparse model 每个存储参数的效率通常低于 dense model，因此它更适合预训练和中高吞吐服务，不一定适合最小设备上的单请求低延迟。

\begin{table}[H]
\centering
\small
\caption{课堂 Q\&A 中的场景化判断。}
\begin{tabularx}{\textwidth}{L{0.28\textwidth}X X}
\toprule
场景 & 可能收益 & 主要风险 \\
\midrule
大规模预训练 & 条件参数扩大容量，data/expert parallel 吞吐高 & all-to-all、存储、checkpoint 与稳定性 \\
中高吞吐服务 & batch 能聚合 expert token，摊薄通信 & 热点 expert、路由漂移与跨机延迟 \\
单请求低延迟/边缘设备 & 可能减少 active FLOPs & 总权重难放下，低 batch 无法形成高效 expert GEMM \\
超长上下文 & MoE 扩大 FFN 容量 & attention/KV 成本可能成为另一独立瓶颈 \\
\bottomrule
\end{tabularx}
\end{table}

\teachervoice{Barret Zoph 的最终判断很克制：如果目标是“在尽可能小的设备上放下最强模型”，dense model 的 per-weight efficiency 可能更好；如果是大量预训练和中高吞吐服务，sparsity 更可能成为大收益。应用场景决定结论。}

\subsection{本章小结}

Switch 的语言结果说明条件参数可以规模化；V-MoE 进一步用 priority routing 分配有限计算。Q\&A 则提醒，存储、batch、互连、上下文长度和每权重效率决定 sparse deployment 是否真正占优。

